\secnumbersection{CONCLUSIONES}

\subsection{SÍNTESIS DEL TRABAJO REALIZADO}

A partir de la problemática de accesibilidad web para personas con discapacidad visual, validada empíricamente mediante la encuesta propia del Capítulo 1, se planteó la idea de un relator de pantalla: un asistente basado en un modelo de lenguaje visual (VLM) que no existía previamente bajo esta forma, diseñado para coexistir con el lector de pantalla del usuario en lugar de reemplazarlo -entregando la respuesta a través del mismo mecanismo de accesibilidad del sistema operativo que este ya utiliza-, en línea con el principio de uso equitativo del diseño universal.

Esa idea se implementó de punta a punta mediante una extensión de navegador, un servidor de gestión y un servidor de inferencia propio (\textit{albert}), con dos modalidades de interacción (texto y voz). Para seleccionar el modelo VLM del sistema se realizó un proceso de \textit{benchmark} comparativo entre 5 candidatos, evaluado mediante un juez LLM (gemma3-27b) diseñado y calibrado especialmente para esta tarea, validado contra un criterio humano mediante el esquema de \textit{LLM-as-a-judge-evaluation} a lo largo de 23 versiones de \textit{prompt}. Este juez se estancó en un 72,5\% de \textit{exact match} sin acercarse al 100\%, techo que se interpreta como el máximo potencial de esta implementación particular; se detectó además que parte de su mejor resultado histórico dependía de un sesgo de contexto en los ejemplos de su \textit{prompt}, lo que obligó a una segunda ronda de calibración.

Finalmente, se validó el sistema mediante un conjunto de métricas técnicas (exactitud de STT, latencia del servidor de gestión, latencia de captura, pertinencia de las respuestas) y mediante entrevistas con 3 personas ciegas que usaron el sistema en tareas reales sobre sitios web. Las entrevistas mostraron una impresión general positiva -la mayoría de las tareas intentadas se completaron o al menos se lograron parcialmente-, con buen desempeño en tareas descriptivas o de extracción de información, aunque con un desempeño más débil en tareas procedimentales o de navegación y en portales que requieren autenticación. Entre los modos de falla identificados se cuentan un sesgo de perspectiva visual en las respuestas del modelo, alucinaciones cuando el elemento consultado no estaba visible en la captura, descripciones superficiales en sitios con estructuras más complejas, y fallas repetidas del mecanismo de \textit{scroll \& stitch} al capturar ventanas emergentes.

Algunos aspectos, además, quedaron inconclusos al cierre de este documento. La validación con usuarios se realizó con solo 3 participantes, por lo que sus resultados se reportan como casos ilustrativos y no como evidencia estadísticamente representativa. Algunas sugerencias reportadas por los propios participantes -como el contexto de conversación- se plantean como trabajo futuro en la siguiente sección.

\subsection{TRABAJO FUTURO}

A partir de lo anterior, se identifican las siguientes líneas de trabajo futuro:

\begin{itemize}
    \item Ampliar la validación con usuarios a una muestra mayor, que permita un análisis más representativo que los casos ilustrativos reportados en este trabajo.
    \item Corregir la falla del mecanismo de \textit{scroll \& stitch} al capturar ventanas emergentes.
    \item Reducir el sesgo de perspectiva visual y mejorar el desempeño del sistema en tareas procedimentales o de navegación, y en portales que requieren autenticación.
    \item Incorporar contexto de conversación, para permitir preguntas de seguimiento sobre un mismo elemento sin repetir la información ya entregada.
    \item Mejorar el \textit{feedback} de la entrada por voz, para que el usuario sepa si su pregunta se está capturando correctamente.
    \item Desarrollar una versión compatible con dispositivos móviles y aplicaciones, además de los navegadores de escritorio.
    \item Explorar arquitecturas de juez LLM con acceso directo a la imagen evaluada, para reducir la dependencia de un juez \textit{text-only} frente a alucinaciones del VLM que sean difíciles de detectar solo a partir del texto.
\end{itemize}